\bookchapter{Chapitre 19 : Dépannage et sauvetage}

\chapterintro{Le smartphone, comme l'homme, tombe parfois « malade ». Ralentissements, fermetures inattendues, écran noir, batterie qui ne charge plus... la plupart des pannes ne sont pas si terribles, vous pouvez tout à fait essayer par vous-même avant de l'envoyer en réparation. Ce chapitre est votre « manuel d'urgence smartphone » — essayez d'abord les méthodes à faible coût, si ça ne marche pas, consultez un médecin professionnel.}

\booksection{I. Le smartphone ralentit, chauffe — la question la plus fréquente}

Honnêtement, après environ un an d'utilisation, le téléphone commence à ralentir, c'est tout à fait normal. La raison n'est pas forcément que le téléphone est cassé, c'est plus souvent « accumulation de poussière » — fichiers inutiles, processus en arrière-plan, espace de stockage qui se remplit progressivement.

\booksubsection{D'abord le diagnostic, ne vous précipitez pas pour réparer :}

\begin{enumerate}
  \item \textbf{Vérifier l'espace de stockage}: allez dans Réglages Stockage pour voir l'espace restant. S'il est inférieur à 10 \%, le système devient très lent car il peine à libérer de l'espace. Mon expérience : maintenez au moins 15 \%-20 \% d'espace libre pour que le téléphone reste fluide.
  \item \textbf{Vérifier le classement de consommation}: allez dans les réglages de batterie pour voir le classement de consommation des dernières 24 heures. Si une application représente la part du lion mais que vous ne l'avez presque pas utilisée, c'est qu'elle tourne en arrière-plan. Désinstallez-la ou limitez son activité en arrière-plan.
  \item \textbf{Rappelez-vous ce que vous avez fait récemment}: est-ce après une mise à jour système que c'est devenu lent ? (après une mise à jour majeure, c'est normal que ce soit lent les premiers jours, le système reconstruit les index en arrière-plan). Ou après avoir installé une nouvelle application ? (essayez de la désinstaller).
  \item \textbf{Observez l'environnement d'utilisation}: le téléphone exposé au soleil, dans une coque épaisse, en chargeant et en jouant, chauffe inévitablement. Les températures élevées provoquent également une baisse de fréquence automatique du CPU, vous sentez le téléphone soudainement ralentir. Laissez-le refroidir puis réessayez.
\end{enumerate}

\textbf{Méthodes de traitement (du simple au complexe)}: redémarrez le téléphone (résout 60 \% des problèmes) nettoyez l'espace de stockage (supprimez les applications et vidéos inutiles) si une application fait des siennes, videz son cache ou désinstallez et réinstallez dernier recours : après sauvegarde des données, réinitialisez aux paramètres d'usine.

\booksection{II. Fermetures inattendues, applications qui ne répondent plus — l'application désobéit}

Une application qui se ferme dès qu'on l'ouvre, ou qui se fige soudainement pendant l'utilisation. Le processus de traitement est comme une consultation médicale, étape par étape :

\begin{enumerate}
  \item \textbf{D'abord fermer et rouvrir}: faites glisser l'application hors de l'arrière-plan, puis rouvrez-la. La plupart des bugs temporaires sont résolus ainsi.
  \item \textbf{Mettre à jour l'application}: allez dans la boutique d'applications pour voir s'il y a une mise à jour. Parfois, c'est parce que la version de l'application est trop ancienne et incompatible avec le système.
  \item \textbf{Vérifier les permissions}: certaines applications nécessitent une certaine permission pour fonctionner normalement, mais vous ne l'avez pas accordée. Par exemple, une application appareil photo qui se ferme dès l'ouverture si la permission caméra n'est pas accordée. Allez dans les réglages pour trouver cette application, accordez toutes les permissions nécessaires.
  \item \textbf{Vider le cache}: allez dans Réglages Gestion des applications trouvez cette application Stockage Vider le cache. Attention, c'est « vider le cache » et non « effacer les données » ! Effacer les données supprimera l'état de connexion et les réglages de l'application, équivalent à une réinitialisation d'usine de l'application.
  \item \textbf{Désinstaller et réinstaller}: dernier recours. Avant de désinstaller, confirmez d'abord si les données de cette application seront perdues — pour WeChat, sauvegardez d'abord l'historique des discussions, pour les jeux, confirmez d'abord si le compte est lié à un numéro de téléphone.
\end{enumerate}

\booksection{III. Ne charge plus ou charge lentement — ce n'est peut-être pas la faute de la batterie}

Ne plus charger est un problème particulièrement anxiogène, mais souvent la cause vient des accessoires, pas du téléphone en panne.

\begin{itemize}
  \item \textbf{D'abord vérifier le port de charge}: avec un cure-dent ou un petit outil non conducteur, nettoyez délicatement la poussière et les peluches du port de charge. Je l'ai vécu moi-même — un fil tombé de la poche s'était coincé dans le port de charge, en le branchant je pensais à un mauvais contact, en le retirant c'était réglé.
  \item \textbf{Essayez un autre câble et chargeur}: le câble s'abîme plus souvent que le chargeur, surtout les câbles souvent enroulés. Essayez un câble et un chargeur dont vous êtes sûr qu'ils fonctionnent, en test croisé, vous identifierez rapidement quel accessoire est en panne.
  \item \textbf{Problème de charge sans fil}: si vous utilisez la charge sans fil, vérifiez si la coque du téléphone n'est pas trop épaisse, ou s'il n'y a pas de plaque métallique derrière la coque (support magnétique, etc.). De plus, placez le téléphone exactement au centre du pad de charge — s'il est décalé, cela peut ne pas charger.
  \item \textbf{Charge lente}: confirmez que la puissance du chargeur est suffisante (l'iPhone nécessite au moins un chargeur PD de 20 W pour la charge rapide), que le câble prend en charge la charge rapide (tous les câbles ne supportent pas plus de 3 A), que le téléphone n'est pas trop chaud (trop chaud réduit automatiquement la puissance pour protéger la batterie).
\end{itemize}

\booksection{IV. Pas de connexion réseau, WiFi se déconnecte sans cesse}

En extérieur, les données mobiles fonctionnent bien, mais en rentrant chez soi, impossible de se connecter au WiFi ; ou le WiFi est connecté mais pas d'Internet. Vérifiez dans cet ordre :

\begin{enumerate}
  \item \textbf{Activer/désactiver le mode avion}: faites glisser le centre de contrôle vers le bas, activez le mode avion, attendez 5 secondes, désactivez-le. Cela force le téléphone à rechercher à nouveau le réseau, très efficace.
  \item \textbf{« Oublier » le WiFi puis se reconnecter}: allez dans les réglages WiFi, appuyez sur le point d'exclamation à côté du nom du réseau (engrenage sur Android), sélectionnez « ignorer ce réseau » ou « oublier », puis saisissez à nouveau le mot de passe pour vous connecter.
  \item \textbf{Redémarrer le routeur}: beaucoup de routeurs domestiques ne sont pas éteints depuis des mois, avec le temps ils deviennent instables. Débranchez l'alimentation, attendez 30 secondes, rebranchez, la plupart des problèmes WiFi sont résolus.
  \item \textbf{Vérifier si le problème vient du téléphone}: les autres appareils se connectent-ils au même WiFi normalement ? Si oui mais que seul votre téléphone ne fonctionne pas, le téléphone a peut-être besoin d'un redémarrage ou d'une réinitialisation des réglages réseau (Réglages Général Transférer ou réinitialiser Réinitialiser les réglages réseau, ne supprime pas vos données, mais efface tous les mots de passe WiFi).
\end{enumerate}

\booksection{V. Problèmes d'écran : tactile insensible, scintillement, images rémanentes}

\begin{itemize}
  \item \textbf{Tactile insensible}: retirez d'abord le film trempé pour voir si c'est sa faute — les films trempés bon marché peuvent effectivement affecter le tactile après un certain temps. L'eau sur l'écran ou l'huile sur les doigts peuvent également provoquer des dérives tactiles. Si ce n'est pas le cas, redémarrez et réessayez.
  \item \textbf{Marques/rémanences sur l'écran (brûlure d'écran)}: l'affichage prolongé de la même image (par exemple, laisser l'écran allumé en scrolant TikTok) laisse une « image rémanente » sur les écrans OLED. Ce n'est pas un écran cassé, c'est temporaire. Éteignez l'écran et laissez-le reposer une demi-heure, généralement ça se rétablit.
  \item \textbf{Scintillement de l'écran}: si le scintillement se produit à faible luminosité, c'est peut-être un problème de modulation PWM, fréquent sur les écrans OLED. Cherchez l'option « dimming DC » ou « anti-scintillement » dans les réglages, l'activation améliore beaucoup les choses. Tous les téléphones n'ont pas cette fonction.
  \item \textbf{Une zone de l'écran complètement noire/allumée}: c'est plus grave, probablement un dommage physique de l'écran, remplacement nécessaire. Les écrans remplacés en réparation non officielle ont généralement des couleurs et une luminosité inférieures à l'original, et affectent l'étanchéité.
\end{itemize}

\booksection{VI. Impossible de s'allumer, bloqué sur le logo — ne paniquez pas}

Cette situation fait le plus battre le cœur, mais restez calme, essayez d'abord ceci :

\begin{itemize}
  \item \textbf{Chargez pendant une demi-heure d'abord}: souvent, c'est simplement que la batterie est complètement épuisée. Branchez le chargeur d'origine, attendez au moins 30 minutes, puis essayez de l'allumer. Certains téléphones, après épuisement complet, nécessitent une certaine charge pour réagir.
  \item \textbf{Redémarrage forcé}: pour les iPhone plein écran, appuyez sur volume haut, volume bas, puis maintenez le bouton latéral d'alimentation jusqu'à ce que le logo Apple apparaisse (peut nécessiter de maintenir 10-15 secondes). Méthode générale Android : maintenez le bouton d'alimentation pendant plus de 10 secondes, jusqu'à la vibration ou l'allumage de l'écran.
  \item \textbf{Bloqué sur l'interface de démarrage avec redémarrages en boucle}: cette situation s'appelle « pomme blanche » (iPhone) ou « blocage sur le logo » (Android). Si le redémarrage forcé ne fonctionne pas, il faut peut-être entrer en mode de récupération et réparer avec un ordinateur. Connectez l'iPhone à l'ordinateur, ouvrez iTunes/Finder, entrez en mode de récupération et sélectionnez « Mettre à jour » (pas « Restaurer »), cela réinstallera le système en préservant les données. Pour Android, entrez en mode Fastboot ou Recovery, les combinaisons de touches varient selon les marques, recherchez votre modèle de téléphone + « entrer en mode récupération ».
  \item \textbf{Si rien de tout cela ne fonctionne}: il ne reste qu'à faire appel à une réparation professionnelle. Si les données ont été sauvegardées, pas trop d'inquiétudes.
\end{itemize}

\booksection{VII. Problèmes de son : pas de son du haut-parleur, son d'appel faible}

\begin{itemize}
  \item \textbf{Pas de son du haut-parleur}: vérifiez d'abord que ce n'est pas le mode silencieux (la bascule de silence latérale de l'iPhone), le mode ne pas déranger, ou le volume réglé au minimum. Ça a l'air bête, mais cela représente 40 \% des cas. Vérifiez ensuite si le Bluetooth n'est pas connecté à un écouteur ou une enceinte.
  \item \textbf{L'interlocuteur ne vous entend pas}: vérifiez si le trou du microphone n'est pas bouché — le petit trou en bas du téléphone se bouche facilement avec la poussière. Nettoyez-le délicatement avec une brosse souple. N'utilisez pas d'aiguille ou de cure-dent pour pousser, vous risquez d'endommager le grill anti-poussière.
  \item \textbf{Son rauque/grésillement}: si c'est dû à de l'eau dans le haut-parleur — n'utilisez pas de sèche-cheveux à air chaud ! Absorbez l'eau en surface avec un essuie-tout, puis laissez sécher naturellement. Sur Internet, on dit que l'application d'évacuation d'eau par ondes sonores a un certain effet (jouer des sons basse fréquence pour faire vibrer l'eau hors du haut-parleur), vous pouvez essayer.
\end{itemize}

\booksection{VIII. Réparation officielle, évitez les piègs}

Si toutes les méthodes ci-dessus ont été essayées sans succès, il faut faire appel à une réparation professionnelle. Mais le marché de la réparation est complexe, je vais parler de mes propres expériences et pièges rencontrés :

\begin{itemize}
  \item \textbf{Trouvez le vrai officiel, pas « ceux qui ressemblent à de l'officiel »}: allez sur le site officiel de la marque (Apple.com.cn, Huawei.com.cn, etc.) pour trouver la liste des « prestataires de services autorisés ». Les boutiques de rue avec un logo Apple mais non autorisées peuvent utiliser des pièces non originales, après réparation votre étanchéité est fichue.
  \item \textbf{Ce qu'il faut faire avant de confier le téléphone}: sauvegardez toutes les données (ne comptez pas sur le fait que les données ne seront pas perdues pendant la réparation), déconnectez tous les comptes importants (identifiant Apple, WeChat, Alipay), désactivez « Localiser mon téléphone ». Certains réparateurs peuvent vous demander le code de déverrouillage, cela pose un risque de vie privée — mieux vaut vous déconnecter d'abord des comptes avant de confier le téléphone.
  \item \textbf{Les pièges du remplacement d'écran}: les écrans non officiels se divisent en « écran pressé original » (verre original remplacé), « écran pressé arrière » (verre non original), « écran assemblé » (tout non original). Les prix varient du simple au triple, mais la qualité d'affichage et la durabilité varient aussi beaucoup. Si le téléphone est encore sous garantie, restez sur l'officiel.
  \item \textbf{Les pièges du remplacement de batterie}: les batteries non originales peuvent surestimer la capacité, et après six mois, elles se dégradent plus que les batteries originales après deux ans. Le prix officiel du remplacement de batterie Apple n'est pas beaucoup plus élevé que le non officiel, il est conseillé de passer par l'officiel.
  \item \textbf{Les pièges de la récupération de données}: après un téléphone tombé dans l'eau/endommagé, si le réparateur dit « la récupération de données coûte plus cher », attention — certains commerçants malhonnêtes exagèrent délibérément l'ampleur des dégâts pour facturer plus. Une récupération de données sérieuse facture après devis, et ne facture pas si la récupération échoue.
\end{itemize}

\begin{tcbwarn}{La première astuce universelle : redémarrez le téléphone}
Quel que soit le problème que vous rencontrez — ralentissement, fermeture inattendue, connexion réseau impossible, son anormal — redémarrez d'abord le téléphone. Cela ressemble à un repoussoir, mais d'après mon expérience, au moins six petits problèmes sur dix sont résolus après un redémarrage. Le téléphone, comme l'ordinateur, accumule avec le temps toutes sortes de choses en mémoire, le redémarrage permet de tout effacer et de recommencer. Donc, la prochaine fois que votre téléphone a un problème, ne vous précipitez pas pour chercher « quoi faire », redémarrez d'abord puis regardez.

\end{tcbwarn}

\begin{tcbpractice}{Mémorisez maintenant : ces opérations de sauvetage doivent être connues par cœur}
\begin{enumerate}
  \item \textbf{Redémarrage forcé}: iPhone plein écran, appuyez sur volume haut volume bas maintenez le bouton d'alimentation jusqu'à l'apparition du logo. Android, maintenez le bouton d'alimentation 10-15 secondes jusqu'à la vibration.
  \item \textbf{Nettoyer le port de charge}: trouvez un petit outil (cure-dent ou épingle d'extraction de carte SIM), nettoyez délicatement la poussière du port de charge.
  \item \textbf{Reconnexion WiFi}: allez dans les réglages WiFi, « oubliez » votre WiFi actuel, saisissez à nouveau le mot de passe pour vous connecter.
  \item \textbf{Activer/désactiver le mode avion}: faites glisser le centre de contrôle vers le bas, activez le mode avion, comptez 5 secondes mentalement, désactivez-le, voyez si le réseau se rétablit.
  \item Capturez les étapes de ces opérations en images, ou écrivez-les dans les notes. En cas de problème, pas besoin de chercher fébrilement, ouvrez simplement les notes.
\end{enumerate}

\end{tcbpractice}
